(1) Freeness: a relation $\sum c_i\vartheta_i = 0$ has interior coefficients $c_0, c_1, c_2\tau^{p_2}, c_3\tau^{p_3}$ against the coframe, all forced to vanish where $\tau > 0$, hence identically. The Gram entries: $\vartheta_2$ carries $\tau^{2p_2}(-\tau^{-2p_2}c_2) = -c_2$, likewise $\vartheta_3$; and $\varrho \in \mathcal{A}$ since $-2p_1 > 0$ pointwise. Corank and radical are read off the boundary values $\operatorname{diag}(1, 0, -c_2, -c_3)$. (2) The coupling construction is a pointwise modification of the Gram by elements of $\mathcal{A}$ and constants; the determinant identity of Lemma \claimref{sheaf.determinant-identity-blow-up-expansion} is algebraic, valid over any commutative ring, with $s = G^{-1}(e,e) = -1/(\varrho c_1)$ here; the boundary signature is Lemma \claimref{sheaf.wall-gram} with $\mu\nu < 0$; the global determinant is the same identity at $\beta \equiv 1$. (3) The diagonal entries have $\tau$-orders $(0, -2p_1, -2p_2, -2p_3)$; each regime's named profile prescribes the target order of every slot, and the resulting equations are linear with the displayed unique solutions. All computational clauses are verified as exact rational-function identities at three rational Kasner points, with the ladder pair on all $64$ boundary states (Section \ref{sec:verification}).
